\section{Related Works}
\label{sec:related}
\paragraph*{GNNs for Physics-Based Simulation}
The application of GNNs to physics-based simulation has first been applied to deformable solids and fluids (both represented by particles)~\cite{sanchez2018graph}. A notable milestone in this field is \flatGNN~\cite{pfaff2020learning}, which enables the general scheme for learning mesh-based simulations. Subsequently, several variants of \flatGNN~have been proposed: such as combining GNNs with Physics-Informed Neural Networks (PINNs)~\cite{gao2022physics}, making long-term predictions by combining the GraphAutoEncoder (GAE) and Transformer~\cite{han2022predicting}, directly predicting the steady-states through the multi-layer readouts~\cite{harsch2021direct}, and accelerating the finer-level simulation by feeding the up-sampled coarser results inferred by GNN~\cite{belbute2020combining}.

\paragraph*{Multi-Scale GNNs}
Multi-scale GNNs (MS-GNNs) have been widely used in general graph-related tasks other than physics~\cite{wu2020comprehensive,mesquita2020rethinking,zhang2019hierarchical}. The GraphUNet\cite{gao2019graph} introduces the UNet structure into GNNs with a trainable scoring module for pooling, and a $2^{\text{nd}}$-powered adjacency enhancement to help conserve the connectivity. Several works have investigated the use of MS-GNNs for physics-based simulations, including the two- and multi-level GNNs\cite{fortunato2022multiscale,liu2021multi}, which rely on manually drawing coarse meshes. Works such as \linoGNN~\cite{lino2021simulating,lino2022towards} rely on spatial proximity to generate multi-level structures. \citet{li2020multipole} adopts multi-level matrix factorization to generate the kernels at coarser levels. \citet{lino2022multi} utilizes Guillard’s coarsening algorithm to build the coarse-level meshes, but only for 2-D triangle elements. Additionally, there are representative works that abandon the original mesh and build connections and hierarchies on point clouds, such as GNS~\cite{sanchez2020learning}, PointNet~\cite{qi2017pointnet}, PointNet++\cite{qi2017pointnet++}, and GeodesicConv\cite{masci2015geodesic}.
